% =========================================================================
% Solución: C1_E01_P04 - Disco Rotante con Punto en Movimiento Radial
% =========================================================================

\begin{sol}[Tema: Cinemática en Coordenadas Cilíndricas - Disco Rotante]

\subsubsection*{Pregunta 1: Velocidad del punto P}

En coordenadas cilíndricas, la posición del punto es:
$$\vec{r} = r(t) \, \hat{e}_r + z \, \hat{e}_z$$

Con $r(t) = r_0 + vt$ (movimiento radial constante) y $\theta(t) = \Omega t$ (rotación a velocidad angular constante).

La velocidad en coordenadas cilíndricas es:
$$\vec{v} = \dot{r} \, \hat{e}_r + r \dot{\theta} \, \hat{e}_\theta + \dot{z} \, \hat{e}_z$$

Sustituyendo nuestras funciones:
\begin{equation}
\boxed{\vec{v} = v \, \hat{e}_r + (r_0 + vt) \Omega \, \hat{e}_\theta}
\end{equation}

Interpretación:
\begin{itemize}
    \item \textbf{Componente radial}: $v_r = v$ (constante, movimiento hacia afuera)
    \item \textbf{Componente transversal}: $v_\theta = (r_0 + vt) \Omega$ (aumenta con el tiempo)
    \item \textbf{Magnitud de velocidad}: $|\vec{v}| = \sqrt{v^2 + [(r_0 + vt) \Omega]^2}$
\end{itemize}

\subsubsection*{Pregunta 2: Aceleración del punto P}

La aceleración en coordenadas cilíndricas es:
$$\vec{a} = (\ddot{r} - r\dot{\theta}^2) \hat{e}_r + (r\ddot{\theta} + 2\dot{r}\dot{\theta}) \hat{e}_\theta + \ddot{z} \, \hat{e}_z$$

Con nuestras funciones:
\begin{itemize}
    \item $\dot{r} = v$ (velocidad radial constante)
    \item $\ddot{r} = 0$ (aceleración radial nula)
    \item $\dot{\theta} = \Omega$ (velocidad angular constante)
    \item $\ddot{\theta} = 0$ (aceleración angular nula)
    \item $r = r_0 + vt$
\end{itemize}

Sustituyendo:
\begin{equation}
\vec{a} = [0 - (r_0 + vt)\Omega^2] \hat{e}_r + [0 + 2v\Omega] \hat{e}_\theta
\end{equation}

\begin{equation}
\boxed{\vec{a} = -(r_0 + vt) \Omega^2 \, \hat{e}_r + 2v\Omega \, \hat{e}_\theta}
\end{equation}

\subsubsection*{Pregunta 3: Componentes de Aceleración}

\paragraph{Componente Radial (Centrípeta)}
$$a_r = -(r_0 + vt) \Omega^2$$

Esta es la aceleración \textbf{centrípeta}: negativa (apunta hacia el centro), proporcional a la distancia radial, y proporcional al cuadrado de la velocidad angular. Su magnitud aumenta linealmente con el tiempo.

\paragraph{Componente Transversal (Coriolis)}
$$a_\theta = 2v\Omega$$

Esta es la aceleración de \textbf{Coriolis}: es constante, proporcional a la velocidad radial $v$ y a la velocidad angular $\Omega$. Representa el cambio de dirección de la velocidad debido a la rotación del sistema.

\paragraph{Componente Axial}
$$a_z = 0$$

No hay movimiento en la dirección axial (perpendicular al disco).

\paragraph{Magnitud Total}
\begin{equation}
|\vec{a}| = \sqrt{[(r_0 + vt) \Omega^2]^2 + (2v\Omega)^2}
\end{equation}

\paragraph{Interpretación Física}

El término centrípeto $-(r_0 + vt) \Omega^2$ es necesario para mantener el punto en una trayectoria circular alrededor del eje mientras se aleja del centro. El término de Coriolis $2v\Omega$ resulta del hecho de que la velocidad del punto en la dirección transversal $v_\theta = (r_0 + vt) \Omega$ cambia continuamente debido al movimiento radial: conforme el punto se aleja del eje, su velocidad tangencial aumenta, requiriendo una aceleración tangencial.

\end{sol}
